\chapter{The find command}

\section{What is it?}

This command allows us to search for files or directories in specified paths with different options to narrow down
the specifics of what we are looking for.

\lstinputlisting[
    language=bash,
    caption={9-bash.sh},
    captionpos=t,
    label={lst:9-bash},
    breaklines=true]
{../9-bash/9-bash.sh}

\section{The find command is very extensive}

There are many options that you can use in order to customize your search without taking into account 
we can also use wild cards in combination. First, lets take a look at the `\underbar{-type}' test expression.

\lstinputlisting[
    language=bash,
    caption={9-bash.sh},
    captionpos=t,
    label={lst:9.1-bash},
    breaklines=true]
{../9-bash/9.1-bash.sh}
\pagebreak
\section{Execute commands if file is found}

There is the `\underbar{-exec}' action expression which allows you to execute any other command if the
find command successfully finds a file (exit code = 0).

\lstinputlisting[
    language=bash,
    caption={9-bash.sh},
    captionpos=t,
    label={lst:9.2-bash},
    breaklines=true]
{../9-bash/9.2-bash.sh}

Alternatively, you could use `$\backslash;$' to terminate (which is mandatory for -exec) if you know there
is only one output from find.

\section{A more practical example of the find command}

say you have a list of images you wish to keep only visible to your user. You'd have to run:

\begin{lstlisting}[
    language=bash, 
    caption={Terminal},
    captionpos=t,
    label={lst:9-terminal}, 
    mathescape=true, 
    breaklines=true]
chmod 600 -R ./Pictures/
\end{lstlisting}

We don't want the file to be executable, that's why we go with 600. There's a problem wiht this approach
whatsoever. The ``Pictures'' directory will also be marked as rw-. This is bad because in order for anyone
to access a directory, it must have the executable option on. This means that we need to apply chmod
600 only to files and ignore directories and subdirectories. We could fix this using:

\begin{lstlisting}[
    language=bash, 
    caption={Terminal},
    captionpos=t,
    label={lst:9.1-terminal}, 
    mathescape=true, 
    breaklines=true]
chmod u+x -R ./Pictures/
\end{lstlisting}

But if we were dealing with files, making them executable introduces vulnerabilities or the chance of
bricking the system. Most likely there is a way to achieve our purpose with the chmod comman itself.
However this is a perfect oportunity to use the find command with the \underbar{-exec} expression.


\lstinputlisting[
    language=bash,
    caption={9-bash.sh},
    captionpos=t,
    label={lst:9.3-bash},
    breaklines=true]
{../9-bash/9.3-bash.sh}

\pagebreak
\section{Using the find command responsibly}

Before using the \underbar{-exec} expression for anything, you should just run the find command as is
and check if the output of it are acutally the files and directories that you wish to remove. Otherwise
you could alter something that you shouldn't have and end up bricking your system